\documentclass[11pt]{article}

\usepackage[T1]{fontenc}
\usepackage[utf8]{inputenc}
\usepackage{lmodern}
\usepackage[a4paper,margin=27mm]{geometry}
\usepackage{microtype}
\usepackage{amsmath,amssymb,amsthm,mathtools}
\usepackage{booktabs}
\usepackage{enumitem}
\usepackage{xcolor}
\usepackage{hyperref}

\definecolor{linkblue}{RGB}{20,67,110}
\hypersetup{colorlinks=true,linkcolor=linkblue,citecolor=linkblue,
  urlcolor=linkblue,
  pdftitle={Oriented Area-Regge Asymptotics of a Native Four-Simplex Vertex},
  pdfauthor={Tobias Croydon-McRae}}

\setlength{\parindent}{1.25em}
\setlength{\parskip}{0.2em}
\setlist{nosep}
\numberwithin{equation}{section}

\newtheorem{theorem}{Theorem}[section]
\newtheorem{proposition}[theorem]{Proposition}
\newtheorem{lemma}[theorem]{Lemma}
\newtheorem{corollary}[theorem]{Corollary}
\newtheorem{conjecture}[theorem]{Conjecture}
\theoremstyle{definition}
\newtheorem{definition}[theorem]{Definition}
\newtheorem{assumption}[theorem]{Standing assumption}
\theoremstyle{remark}
\newtheorem{remark}[theorem]{Remark}

\newcommand{\SUtwo}{SU(2)}
\newcommand{\SpinFour}{\mathrm{Spin}(4)}
\newcommand{\Dom}{\mathcal D}
\renewcommand{\Re}{\operatorname{Re}}
\renewcommand{\Im}{\operatorname{Im}}

\title{\textbf{Oriented Area-Regge Asymptotics\\of a Native Four-Simplex Vertex}\\[0.5em]
\large A single geometric saddle on nondegenerate coherent data}
\author{Tobias Croydon-McRae\\\small Independent researcher}
\date{Preprint v0.1 --- 19 September 2026}

\begin{document}
\maketitle

\begin{abstract}
We study an orientation-resolved Riemannian four-simplex vertex, W2, built from
a finite relational spin model. Orientation is fixed by restricting the history
integral to the domain on which a native $\SpinFour$-invariant pseudoscalar
$\Omega$ has fixed sign. The native positive volume insertion is
$|Q|^{1/2}$ on each tetrahedron. No per-face sign variable, Barrett--Crane
projection, or additional Einstein--Hilbert boundary projector enters the
vertex analysed here.

For boundary data given by Livine--Speziale coherent intertwiners of a
nondegenerate four-simplex, with spins $j_f=\lambda k_f+O(1)$, we show, subject
to the one hypothesis recorded below: (i) the
integrand is exponentially small except near the four pairings of the two chiral
gluing solutions; (ii) exactly one of these, a geometric saddle, lies in the
interior of the orientation domain, and every configuration on the domain
boundary is either exponentially suppressed or rank one, where the measure
vanishes cubically and the contribution is subleading; (iii) the wrong-sign
volume weight of the boundary state is $O(\lambda^{-\infty})$, by an elementary
ladder-operator estimate. The real part of the coherent-action Hessian is proved exactly to be negative
definite modulo gauge on both the regular and the Dittrich--Kogios simplex; this
also makes the full complex gauge-fixed Hessian nonsingular. The vertex is
therefore asymptotic to a single oriented exponential
$e^{\,i\sigma\,2\sum_f j_f\Theta_f}$, the oriented area-Regge phase.

We do \emph{not} obtain a continuum statement. An earlier version inferred, from
gluing, the linearised graviton dynamics that Dittrich and Kogios derive from the
area-Regge action. That inference is withdrawn: contracting coherent intertwiners
across a shared tetrahedron forces the two shapes to agree, whereas the
Dittrich--Kogios configuration space relies on the shape mismatch such gluing
suppresses (Section~\ref{sec:r4}).

The result is Riemannian and leading-order, holds for coherent boundary data of
nondegenerate simplices, and is a vertex statement: it does not control the bulk
spin sum. For the model's own selected tetrahedral state the corresponding
sign-concentration statement is reduced exactly to a spectral conjecture that we
verify for $j\le20$ but do not prove.

No unproved Hessian hypothesis remains at the stated regular/DK scope. The
real-part quadratic form is an explicit weighted sum of squares, and exact
bivector determinants show that its kernel is precisely the common $\SUtwo$
gauge in each chirality. Negative definiteness of the real part then implies
invertibility of the complex gauge-fixed Hessian. The uniform norm asymptotic
that an earlier version assumed is likewise proved as Lemma~\ref{lem:unif}.
\end{abstract}

\tableofcontents

% ===========================================================================
\section{Introduction}\label{sec:intro}

\paragraph{The claim, and its boundary.}
We show that one specific four-simplex vertex, derived inside a finite
relational spin model, is dominated at large spin by a single nondegenerate
geometric saddle whose phase is the oriented area-Regge action. It is natural to
hope that this places the vertex at the start of the chain
\[
\text{W2}\;\longrightarrow\;\text{oriented area Regge}\;\longrightarrow\;
\text{linearised Riemannian graviton dynamics},
\]
whose second arrow is a theorem of Dittrich and Kogios~\cite{DittrichKogios2023}.
\textbf{We do not establish the first arrow}: gluing coherent vertices forces
matched tetrahedral shapes, which is precisely the freedom the Dittrich--Kogios
configuration space needs (Section~\ref{sec:r4}). We do not claim to derive
general relativity.
The statement is Riemannian, semiclassical, leading-order and linearised; it holds
for coherent boundary data of nondegenerate simplices; it controls each vertex but
not the bulk spin sum of glued histories; and it says nothing about Lorentzian
signature or nonlinear curvature. These limits are collected in
Section~\ref{sec:claims}.

\paragraph{What is standard.}
Most of the chain is not new, and we say so where it appears: coherent
intertwiners and their norm asymptotics~\cite{LivineSpeziale2007}; the two
parity-related solutions of the chiral gluing equations for nondegenerate
Regge-like data~\cite{BarrettEtAl2009}; the Schläfli identity~\cite{Regge1961};
the use of a projector to select the Einstein--Hilbert sector in Engle's
proper vertex~\cite{Engle2011}; and the area-Regge
refinement theorem~\cite{DittrichKogios2023}.

Two further ingredients are standard and we claim nothing for them. The algebra
generated by the pair operators and $Q$ on a four-valent intertwiner is a
realisation of the \emph{Racah algebra} $R(3)$, whose generators are the
intermediate Casimirs and which is precisely the commutant of the diagonal
action~\cite{Racah2019}; the factorisation of Proposition~\ref{prop:factor} is a
one-line consequence of its relations, and the algebraic structure of the
four-valent volume operator has been studied at
length~\cite{BaezBarrett1999,Schliemann2013}. Likewise the moment bound of
Lemma~\ref{lem:ladder} is ordinary spin-coherent-state estimation; its only
virtue is that it replaces a symbol calculus at the spectral cusp.

\paragraph{What is specific to this construction.}
(a) Orientation is selected by a domain restriction on the integration variables,
through a pseudoscalar $\Omega$ built from the native history variables, rather
than by a projector on the boundary state or a per-face sign. (b) The volume
insertions and the cofactor $|\Omega_c|$ give the native W2 weight. Its geometric
cofactor identity and cubic spacetime-volume power are exact, with no fitted
exponent; an exact foliation-free transcription of the full continuum measure
of~\cite{ShiraziEngle2014} is not needed for the theorem and is not claimed here. We have not
located (a) or (b) in the literature; that is a search result, not a novelty
theorem, and the programme that produced this paper has repeatedly found
native-looking structures to be standard. The contribution we do claim is the
assembly: a single-saddle asymptotic for this vertex in which boundary control
and sign concentration are proved rather than assumed.

\paragraph{On method.}
Every load-bearing numerical statement was frozen with its log, reproduced by an
implementation sharing no code, and submitted to hostile review by reviewers from
more than one model provider, including context-free sessions. Every search for
the absence of something carries a positive control that must first find a known
instance; a search whose control fails issues no verdict. The review trail,
including the findings that changed this text, is in
Appendix~\ref{app:repro}.

\paragraph{AI collaboration.}
The derivations, code, numerical experiments and drafting of this paper were
carried out in collaboration with large language models, from more than one
provider, working under the direction of the author. The author set the
questions, decided which lines to pursue, adjudicated the results, and is
responsible for the content; no model is or could be an author, since authorship
carries accountability for the claims and that rests with a person.

We state this because the method is load-bearing rather than incidental. Every
load-bearing numerical claim was reproduced by an independent implementation
sharing no code path; every search for the absence of something carried a
positive control that had to recover a known instance before any verdict was
issued; and each stage was submitted to adversarial review by models from
different providers, instructed to falsify rather than to improve. Several
results in this paper exist in their present form only because such a review
falsified an earlier version, and those corrections are recorded rather than
quietly absorbed.


% ===========================================================================
\section{The native selector}\label{sec:selector}

On the intertwiner space $\mathcal K$ of a four-valent node let $A,B,C$ be the
three pair operators, trace-centred to $a,b,c$ with $a+b+c=0$, and put
\[
Q=i[a,b]=\mathbf J_1\cdot(\mathbf J_2\times\mathbf J_3),\qquad G=a^2+b^2+c^2 .
\]

These generate a realisation of the Racah algebra $R(3)$: its generators are the
intermediate Casimirs $C^{(12)},C^{(23)}$, which are the pair operators up to a
shift, and the algebra is the commutant of the diagonal action~\cite{Racah2019}.
The following identity is an elementary consequence of the two standard relations
$a+b+c=0$ and $[a,b]=[b,c]=[c,a]$. We record it because it is what makes the
selector positive and its ground state sharply defined, not because it is new.

\begin{proposition}[Exact factorisation]\label{prop:factor}
With $\omega=e^{2\pi i/3}$ and $Z=a+\omega b+\omega^2c$,
\[
ZZ^\dagger=\tfrac32\bigl(G-\sqrt3\,Q\bigr),\qquad
Z^\dagger Z=\tfrac32\bigl(G+\sqrt3\,Q\bigr).
\]
Hence $H_\sigma:=G-\sigma\sqrt3\,Q\ge0$, $H_+=\tfrac23ZZ^\dagger$,
$H_-=\tfrac23Z^\dagger Z$, and
\begin{equation}\label{eq:QZ}
Q=\frac{Z^\dagger Z-ZZ^\dagger}{3\sqrt3}.
\end{equation}
\end{proposition}

\begin{proof}
Direct expansion, using $\omega+\omega^2=-1$, $\omega-\omega^2=i\sqrt3$,
$(a+b+c)^2=0$ (which gives $\{a,b\}+\{b,c\}+\{c,a\}=-G$), and
$[a,b]=[b,c]=[c,a]=-iQ$ (a consequence of $a+b+c=0$). The symmetric part of
$ZZ^\dagger$ is $G+\tfrac12G$ and the antisymmetric part is $-\tfrac32\sqrt3\,Q$;
reversing the order flips the sign of the latter.
\end{proof}

The identity holds for arbitrary external spins. Numerically it was verified for
equal spins $j=\tfrac12,\dots,6$ (residuals $\le2.1\times10^{-12}$), and
\eqref{eq:QZ} for $j\le20$ (residuals $\le3.6\times10^{-15}$).

The ground state of $H_+$ is the model's own tetrahedral state. It enters only
in Section~\ref{sec:selected}, which treats it separately from the main theorem.

% ===========================================================================
\section{The oriented vertex W2}\label{sec:w2}

Let $X_a^\pm\in\SUtwo$, $a=0,\dots,4$, be the doubled history variables and
$N_a=X_a^-(X_a^+)^{-1}$, read as unit vectors in $\mathbb R^4$.

\begin{definition}[Vertex]\label{def:w2}
Put $\Omega_a=(-1)^a\det(N_0,\dots,\widehat{N_a},\dots,N_4)$ and
$\Dom_\sigma=\{X:\operatorname{sgn}\Omega_a(X)=\sigma\ \forall a\}$. For a boundary
state $\Psi$ on $\bigotimes_e\mathcal K_e^+\otimes\mathcal K_e^-$,
\[
\mathcal A_\sigma[\Psi]
=\int_{\Dom_\sigma}\!dX\;|\Omega_c(X)|\;
\mathcal I_X\!\Bigl[\bigl(\textstyle\bigotimes_e\widehat T_e\bigr)\Psi\Bigr],
\qquad \widehat T_e\propto|Q_e|^{1/2}.
\]
where $\mathcal I_X$ is the ordinary doubled $K_5$ matrix element with
$U^\pm_{ab}=D^{j_{ab}}\bigl((X_a^\pm)^{-1}X_b^\pm\bigr)$.
\end{definition}

The following are proved in~\cite{W2vertex}: $\Omega_a$, $\Dom_\sigma$ and
$|\Omega_c|$ are $\SpinFour$ invariant; $\mathcal A_\sigma$ is linear in $\Psi$;
and the exchange $X^+\leftrightarrow X^-$ maps
$\Dom_+\leftrightarrow\Dom_-$. The native volume/cofactor weight is fixed by
the model. Its geometric core can be proved directly: if $T_a$ is the volume of
the facet opposite vertex $a$, $V_4$ the four-volume and $\lambda_a$ the
barycentric coordinate, then
$\nabla\lambda_a=-(T_a/4V_4)N_a$ and
$|\det(\nabla\lambda_b)_{b\ne a}|=(4!V_4)^{-1}$. Hence
\[
 |\Omega_a|
 =\frac{32}{3}\frac{V_4^3T_a}{\prod_bT_b}.
\]
Together with
$T_e^2=(2/9)|\mathbf E_1\!\cdot(\mathbf E_2\times\mathbf E_3)|$
and
$i[\mathbf J_1\!\cdot\!\mathbf J_2,\mathbf J_2\!\cdot\!\mathbf J_3]
=\mathbf J_1\!\cdot\!(\mathbf J_2\times\mathbf J_3)$,
this gives the exact native cubic-volume identity recorded in the companion.
The stronger claim that this is a unique foliation-free discretisation of the
full continuum measure is not used below.

% ===========================================================================
\section{Coherent boundary data and the doubled action}\label{sec:action}

\begin{assumption}\label{ass:data}
Fix a nondegenerate Riemannian four-simplex with face areas $k_f>0$ and outward
face normals $n_{ab}\in S^2$ in each tetrahedral frame. Let $\lambda\to\infty$
along a sequence with $j_f=\lambda k_f+O(1)$ admissible (half-integral and
satisfying the coupling conditions). For each tetrahedron let
$\Phi_e=\bigotimes_{f\ni e}|j_f,n_{ef}\rangle$ and
$\iota_e=\Pi\Phi_e/\|\Pi\Phi_e\|$, $\Pi=\int_{\SUtwo}dg\,U(g)$, be the
Livine--Speziale coherent intertwiner~\cite{LivineSpeziale2007}. The boundary
state is $\Psi=\bigotimes_e\iota_e\otimes\bar\iota_e$.
\end{assumption}

For the Dittrich--Kogios (DK) simplex we take the vertices
\[
(0,0,0,0),\ (1,0,0,0),\ (1,1,0,0),\ (1,1,1,0),\ \tfrac12(1,1,1,1),
\]
whose ten face areas stand in the ratios
$\tfrac{\sqrt3}{4}:\tfrac12:\tfrac{\sqrt2}{2}$ with multiplicities $4:4:2$. We
identify the simplex by its coordinates, not by its areas, since ten areas do not
in general determine a four-simplex. The area ratios are irrational, which is why
the spin sequence is $\lambda k_f+O(1)$ and not $\lambda k_f$.

With $z^\chi_{ab}(X)=\langle Jn_{ab}|(X_a^\chi)^{-1}X_b^\chi|n_{ba}\rangle$, the
coherent component of the integrand is
$\prod_{a<b}(z^+_{ab})^{2j_{ab}}\,\overline{(z^-_{ab})^{2j_{ab}}}
=\exp\bigl[\lambda S(X)+O(1)\bigr]$ with the doubled action
\begin{equation}\label{eq:action}
S(X)=S_+(X^+)+\overline{S_-(X^-)},\qquad
S_\chi=\sum_{a<b}2k_{ab}\log z^\chi_{ab}.
\end{equation}
The action separates additively in the two chiralities, and since
$|z^\chi_{ab}|\le1$, $\Re S=\Re S_++\Re S_-\le0$.

% ===========================================================================
\section{The critical set}\label{sec:critical}

\begin{proposition}\label{prop:crit}
$\Re S(X)=0$ if and only if both chiral families satisfy the gluing equations
$R(X_a^\chi)n_{ab}=-R(X_b^\chi)n_{ba}$. For nondegenerate Regge-like data each
chirality has exactly two solutions up to gauge, related by
parity~\cite{BarrettEtAl2009}. The zero set of $\Re S$ therefore consists of four symmetry classes. (Not four
orbits: each $X_a^\chi$ admits an independent central lift
$X_a^\chi\mapsto\epsilon_a^\chi X_a^\chi$, $\epsilon_a^\chi=\pm1$, which leaves
$\Re S$ unchanged and, when $\epsilon_a^-=\epsilon_a^+$, leaves every $N_a$
unchanged as well, so each class contains finitely many disconnected lift copies
inside the same $\Dom_\sigma$. That changes a finite multiplicity, not the
asymptotic power or phase; the cutoffs in Theorem~\ref{thm:boundary} are taken as
unions over the lift copies.) The classes are:
\begin{enumerate}
\item two \emph{cross} pairings, which reconstruct the two opposite orientations
of the nondegenerate simplex, have all $\Omega_a\neq0$ with a common sign, and lie
one in $\Dom_+$ and one in $\Dom_-$;
\item two \emph{same/same} pairings, at which $N_a=Y_-Y_+^{-1}$ is independent of
$a$ (rank one), so every $\Omega_a=0$.
\end{enumerate}
\end{proposition}

\begin{proof}
$\Re S=0$ forces $|z^\chi_{ab}|=1$ for every face and chirality, which is
saturation of Cauchy--Schwarz, i.e.\ the gluing equation. Both chiralities carry
the same boundary normals, so they have the same two solutions $\{X_a\},\{X'_a\}$.
For the pairing $(X,X)$, $N_a=Y_-X_aX_a^{-1}Y_+^{-1}=Y_-Y_+^{-1}$ for every $a$;
five parallel vectors make every $4\times4$ minor vanish. The cross statement is
the reconstruction theorem of~\cite{BarrettEtAl2009}, and the sign of $\Omega$ is
that of the reconstructed four-orientation~\cite{W2vertex}.
\end{proof}

As an independent check, 4000 random starts on the DK gluing equations found
exactly two gauge-inequivalent chiral solutions.

\paragraph{Phase at the geometric saddle.}
At the cross pairing in $\Dom_\sigma$ the chiral phases are opposite, so
$\Im S=2\sigma\sum_fk_f\Theta_f$ and the oscillatory factor is
$\exp[i\sigma\,2\sum_fj_f\Theta_f]$. We keep the area label $A_f=2j_f$, which is
an integer, so the $2\pi$-periodicity argument of Proposition~\ref{prop:r3}
applies unchanged. The balanced label $2j_f+1$ differs by
$e^{-i\sigma\sum_f\Theta_f}$; that factor is $O(\lambda^0)$ at fixed geometry but is
\emph{not} independent of the bulk once areas are varied, so it may not be
absorbed into a boundary convention in Section~\ref{sec:r3}.

% ===========================================================================
\section{Sign concentration of the volume operator}\label{sec:s5}

Fix a tetrahedron $e$ with classical oriented volume $q_e\neq0$; take $q_e>0$,
the other case being conjugate, and put $\mu_\lambda=\lambda^3q_e$. For a state
$v$ write $p_{\le0}(v)=\|P_{(-\infty,0]}(Q)v\|^2/\|v\|^2$. The zero eigenspace is
included on purpose: at integer spin $Q$ has a zero eigenvalue, so a bound on
$P_{(-\infty,0)}$ alone would not give $p_+\to1$.

\begin{lemma}[Ladder estimate]\label{lem:ladder}
Let $V_k\subset\bigotimes_fV_{j_f}$ be spanned by vectors obtained from $\Phi_e$
by at most $k$ lowering operators, each face in its own frame. Then
$(Q-\mu_\lambda)V_k\subset V_{k+3}$ and
$\|(Q-\mu_\lambda)|_{V_k}\|\le C_k\lambda^{5/2}$, with $C_k$ independent of
$\lambda$.
\end{lemma}

\begin{proof}
Write $\mathbf J_f=(j_f-N_f)\,\mathbf n_f+\mathbf J_f^\perp$ in the frame of face
$f$, where $N_f$ counts lowerings and $\mathbf J_f^\perp$ is built from $J_f^\pm$.
On $V_k$: $0\le N_f\le k$, $\|J_f^-\|\le\sqrt{(k+1)\,2j_f}$ and
$\|J_f^+\|\le\sqrt{k(2j_f+1)}$, both $O(\lambda^{1/2})$. The three faces act on
different tensor factors, so the triple product has no ordering ambiguity.
Expanding $Q=\mathbf J_1\cdot(\mathbf J_2\times\mathbf J_3)$: the term with three
longitudinal factors is $\prod_f(j_f-N_f)\,q=\mu_\lambda+O(\lambda^2)$ on $V_k$,
using $j_f=\lambda k_f+O(1)$; terms with one transverse factor are
$O(\lambda^{5/2})$; with two, $O(\lambda^2)$; with three, $O(\lambda^{3/2})$. Each
term raises the lowering count by at most three.
\end{proof}

\begin{lemma}[Uniform coherent-intertwiner norm]\label{lem:unif}
Let $n_f$, $f=1,\dots,4$, be fixed normals with weights $k_f>0$ satisfying
closure $\sum_fk_fn_f=0$ and spanning $\mathbb R^3$, and let
$j_f=\lambda k_f+\delta_f$ be admissible with $|\delta_f|\le D$. Then there are
$c>0$ and $\lambda_0$, depending only on $(n_f,k_f,D)$, such that
\[
\|\Pi\Phi_\lambda\|^2\;\ge\;c\,\lambda^{-3/2}
\qquad\text{for all }\lambda\ge\lambda_0 .
\]
\end{lemma}

\begin{proof}
Since $\Pi$ is an orthogonal projector,
$\|\Pi\Phi\|^2=\langle\Phi|\Pi|\Phi\rangle=\int_{\SUtwo}dg\prod_fz_f(g)^{2j_f}$
with $z_f(g)=\langle n_f|g|n_f\rangle$ the spin-$\tfrac12$ overlap, $|z_f|\le1$.
Write the integrand as $\exp[\lambda S_0(g)+T_\delta(g)]$ with
\[
S_0=\sum_f2k_f\log z_f,\qquad T_\delta=\sum_f2\delta_f\log z_f .
\]
$\Re S_0\le0$, with equality exactly at $g=\pm1$; at $g=1$ the first variation is
$2i\sum_fk_fn_f=0$ by closure, and the Hessian of $\Re S_0$ is negative definite
because the $n_f$ span, which is the nondegenerate closure saddle
of~\cite{LivineSpeziale2007}. The perturbation $T_\delta$ is holomorphic near
$g=1$ with $T_\delta(1)=0$, and it and its first two derivatives are bounded by
$CD$ uniformly in $\delta$.

Solving $\lambda\,dS_0+dT_\delta=0$ moves the critical point by
$u_\delta=O(D/\lambda)$, and the exponent there is $O(D^2/\lambda)$. Ordinary
three-dimensional stationary phase with a nondegenerate complex Hessian then
gives
\[
\int_{|g-1|<\rho}dg\,e^{\lambda S_0+T_\delta}
=\Bigl(\frac{2\pi}{\lambda}\Bigr)^{3/2}
\frac{e^{T_\delta(1)}}{\sqrt{\det(-S_0'')}}\bigl(1+O(D^2/\lambda)\bigr),
\]
and the complement of the two balls contributes $O(e^{-\eta\lambda})$ by
compactness. The point $g=-1$ contributes the same expression times
$(-1)^{\sum_f2j_f}$, which is $+1$ precisely when the spins are admissible, so
the two add rather than cancel. The prefactor depends continuously on
$\delta$ over the compact set $|\delta_f|\le D$ and is positive there, so its
infimum $c>0$ is attained, which is the claim.
\end{proof}

\begin{remark}
Two earlier versions of this passage were wrong. The first argued that no uniform
constant was needed because the offsets take finitely many values at each
$\lambda$; finiteness at each $\lambda$ gives no single $c$ along an infinite
sequence. The second assumed the uniform statement outright. The point that makes
the proof work is that the offsets enter only through the bounded holomorphic
perturbation $T_\delta$, which shifts the saddle by $O(\lambda^{-1})$ and changes
the prefactor continuously; in particular the $j$-weighted closure defect left by
bounded offsets is $O(1)$ against a Hessian of size $O(\lambda)$. The degenerate
collinear case, which scales as $\lambda^{-1}$, is excluded by the spanning
hypothesis.
\end{remark}

\begin{lemma}[Coherent sign concentration]\label{lem:conc}
For each fixed $n$,
$\langle(Q-\mu_\lambda)^{2n}\rangle_{\iota_e}=O(\lambda^{5n+3/2})$, and hence
$p_{\le0}(\iota_e)=O(\lambda^{-\infty})$.
\end{lemma}

\begin{proof}
$Q$ commutes with the diagonal $\SUtwo$ action, hence with the orthogonal
projector $\Pi$. Therefore
\[
\langle(Q-\mu)^{2n}\rangle_{\iota}
=\frac{\|\Pi(Q-\mu)^n\Phi\|^2}{\|\Pi\Phi\|^2}
\le\frac{\|(Q-\mu)^n\Phi\|^2}{\|\Pi\Phi\|^2}
\le\frac{\bigl(\prod_{i<n}C_{3i}\bigr)^2\lambda^{5n}}{\|\Pi\Phi\|^2}
\]
by Lemma~\ref{lem:ladder}, together with the uniform norm bound of
Lemma~\ref{lem:unif} below. On the range of $P_{(-\infty,0]}(Q)$ one has
$|Q-\mu|\ge\mu$, so by the spectral theorem
$p_{\le0}\le\langle(Q-\mu)^{2n}\rangle/\mu^{2n}=O(\lambda^{3/2-n})$ for every $n$.
\end{proof}

The only input beyond elementary spin algebra is the norm lower bound, the
standard stationary-phase statement for coherent intertwiners. The projection
costs a factor $\lambda^{3/2}$, absorbed because $n$ is arbitrary. No uniformity
of the constants in $n$ is used or needed: $O(\lambda^{-\infty})$ means that for
each $N$ some constant $C_N$ bounds $\lambda^Np_{\le0}$.

\paragraph{Reduction to coherent components.}
The volume weight $\widehat T_e$ is not a polynomial, so
$(\bigotimes\widehat T_e)\Psi$ is not a coherent state. Taylor-expand $|q|^{1/2}$
about $\mu_\lambda$ to order $M$; the remainder is bounded spectrally by
$C_M|q-\mu|^M\mu^{1/2-M}$, and its contribution is $O(\lambda^{3/2-M/2})$ relative
to the leading term by Lemma~\ref{lem:conc}. The Taylor polynomial applied to
$\iota_e$ gives finitely many vectors in $\Pi V_{3M}$, whose matrix elements are
$e^{\lambda S(X)}$ times polynomials in $X$ of bounded degree. Hence, for every
$N$,
\begin{equation}\label{eq:reduce}
\mathcal I_X\bigl[(\otimes\widehat T_e)\Psi\bigr]
=\sum_{k\le K_N}a_k(\lambda)\,P_k(X)\,e^{\lambda S(X)}+R_N(X),\qquad
\sup_X|R_N|\le C_N\lambda^{-N}|a_0(\lambda)|,
\end{equation}
with $a_0\neq0$ the leading coefficient.

\begin{lemma}[Coefficient ratios]\label{lem:coeff}
For every fixed $k\ge1$ and every $\varepsilon>0$,
\[
\Bigl|\frac{a_k}{a_0}\Bigr|=O\bigl(\lambda^{-k/2+k\varepsilon}\bigr).
\]
In particular $a_1/a_0=O(\lambda^{-1/2+\varepsilon})$ and every higher term is
smaller, so the $k=0$ term dominates \eqref{eq:reduce}.
\end{lemma}

\begin{proof}
Put $X=(Q-\mu)/\mu$. Since $\mu\asymp\lambda^3$ and $\|Q\|=O(\lambda^3)$ we have
$\|X\|=O(1)$, and Lemma~\ref{lem:conc} gives
$\langle|X|^{2n}\rangle=O(\lambda^{3/2-n})$ for every fixed $n$. Fix
$\alpha<1/2$ and split the spectral measure of $X$ at $|X|=\lambda^{-\alpha}$. On
the small part, $\|X^k\|\le\lambda^{-\alpha k}$. On the large part Markov gives
\[
\Pr\bigl(|X|>\lambda^{-\alpha}\bigr)
\le\lambda^{2\alpha n}\langle|X|^{2n}\rangle
=O\bigl(\lambda^{3/2-(1-2\alpha)n}\bigr),
\]
which is smaller than any fixed power once $n$ is chosen large enough, while
$\|X\|^k=O(1)$ there. Hence $\|X^k\iota\|=O(\lambda^{-\alpha k})$, and since the
$k$-th Taylor term is $c_k\mu^{1/2}X^k$ against $a_0=\mu^{1/2}$, the stated bound
follows with $\varepsilon=1/2-\alpha$.
\end{proof}

\begin{remark}
An earlier version of this lemma bounded the $k=1$ term by the scalar first
moment $\langle Q\rangle-\mu=O(\lambda^2)$. That is invalid: the Taylor term
involves the vector $\mu^{-1/2}(Q-\mu)\iota$, and a small expectation value does
not bound a norm; used naively Lemma~\ref{lem:conc} would give only
$O(\lambda^{1/4})$, which does not even tend to zero. The interpolation above
avoids the first moment entirely. We are grateful to the reviewer who supplied
it.
\end{remark}

% ===========================================================================
\section{The domain boundary}\label{sec:boundary}

The domain $\Dom_\sigma$ has a boundary, the union of the five cofactor zero
loci. Interiority of the geometric saddle validates its local expansion but does
not by itself bound the rest of the integral. Review~\cite{PR355} pointed out that
a cofactor-order lemma prices certain partly parallel boundary strata at
$\lambda^{-7/2}$, ahead of the geometric saddle. That lemma concerns the
Barrett--Steele mixed-sign families~\cite{BarrettSteele2003}, which exist only
when the coherent sign is a free per-face variable. W2 has no such variable.

\begin{theorem}[Boundary control]\label{thm:boundary}
Under Assumption~\ref{ass:data}, the contribution to $\mathcal A_\sigma$ from any
region of $\partial\Dom_\sigma$ whose closure stays a positive distance from the
rank-one orbits is $O(e^{-\eta\lambda})$ for some $\eta>0$; deleting those orbits
while allowing accumulation would not give a uniform $\eta$. The contribution
from a neighbourhood of the rank-one orbits is smaller than that of the geometric
saddle by a factor $O(\lambda^{-3/2})$. The real-Hessian input
(Proposition~\ref{prop:redef}) is established for the regular and DK data only,
so the second statement carries that restriction.
\end{theorem}

\begin{proof}
Use \eqref{eq:reduce} and a partition of unity
$1=\chi_{\mathrm{geom}}+\chi_{\mathrm{rk1}}+\chi_{\mathrm{rest}}$ on the compact
group manifold. On $\operatorname{supp}\chi_{\mathrm{rest}}$, which avoids the
four classes of Proposition~\ref{prop:crit} and their central lifts, the
continuous function $\Re S$ is
strictly negative, hence $\le-\eta$ by compactness, and every term of
\eqref{eq:reduce} is $O(\lambda^me^{-\eta\lambda})$. The geometric orbit in
$\Dom_\sigma$ is interior, since its cofactors are nonzero, so
$\chi_{\mathrm{geom}}$ may be supported inside $\Dom_\sigma$; the conjugate cross
orbit lies in $\Dom_{-\sigma}$ and does not meet $\overline{\Dom_\sigma}$.

For the rank-one orbits we bound the modulus of the integrand, which requires no
knowledge of the domain shape or of the phase. A rank-one orbit is a pairing
$(X,X)$ of one chiral solution with itself, so by \eqref{eq:action} and
Proposition~\ref{prop:redef} the real part of $S$ has a nondegenerate maximum
there modulo gauge:
$\Re S\le-c\,(|u_+|^2+|u_-|^2)$ in gauge-fixed transverse coordinates $u_\pm$,
$c>0$. All five normals then lie within $O(|u_+|+|u_-|)$ of a common vector, and
each cofactor is a $4\times4$ determinant of four such vectors, so
$|\Omega_c|\le C(|u_+|+|u_-|)^3$. With $D=24$ the gauge-fixed dimension,
\[
\Bigl|\int\chi_{\mathrm{rk1}}\Bigr|
\le C'\!\int_{\mathbb R^{D}}(|u|)^3e^{-c\lambda|u|^2}\,du
\;=\;O(\lambda^{-D/2-3/2}),
\]
while the geometric saddle contributes $c_0\lambda^{-D/2}$ with $c_0\neq0$ by
Proposition~\ref{prop:hess} and $|\Omega_c|>0$ there. The polynomials $P_k$ of
\eqref{eq:reduce} are bounded on the compact group, and their coefficients obey
Lemma~\ref{lem:coeff}, so the finite sum is dominated by its $k=0$ term. The
ratio is therefore $O(\lambda^{-3/2})$, independently of the radial count
of~\cite{W2vertex}, which gives the weaker $\lambda^{-5}$ against
$\lambda^{-9/2}$ and is quoted only as a cross-check from a source that is not
yet public.
\end{proof}

The proof uses only Proposition~\ref{prop:crit}. As a falsification test we
sampled $\sup\Re S$ on the two partly parallel strata named in review, holding
the distinct normals at least $r$ apart (Table~\ref{tab:strata}). A positive
control at $r=0$ must recover $\Re S=0$ at the rank-one corner, and does.

\begin{table}[h]
\centering
\begin{tabular}{lcc}
\toprule
& two-pair, $N_1{=}N_2$, $N_3{=}N_4$ & four-parallel, $N_1{=}\dots{=}N_4$\\
\midrule
control, $r=0$ & $-2.0\times10^{-10}$ & $-2.4\times10^{-11}$\\
$\sup\Re S/\sum_f2k_f$, $r\in[0.05,0.8]$ & $-0.040\,r^2$ & $-0.019\,r^2$\\
\bottomrule
\end{tabular}
\caption{DK data. The quadratic coefficient is constant to two figures over the
whole range; $\Re S$ vanishes on these strata only at the rank-one corner.}
\label{tab:strata}
\end{table}

% ===========================================================================
\section{Nondegeneracy of the geometric saddle}\label{sec:s4}

By \eqref{eq:action} the Hessian of $S$ at any doubled geometric saddle is
block diagonal,
$\operatorname{diag}(H_+,\overline{H_-})$: the exponent contains no term
depending jointly on $X^+$ and $X^-$, while $|\Omega_c|$ is an $O(1)$
prefactor. The required nondegeneracy can be established without a numerical
determinant.

\begin{proposition}[Exact real-part coercivity]\label{prop:redef}
At either chiral gluing solution for the regular or the DK simplex, write
$X_a=X_a^{(0)}\exp(i\,\xi_a\!\cdot\!\sigma/2)$ and let $m_{ab}$ be the coherent
ray on the edge $ab$. The quadratic term in the real part of the chiral action
is, up to the fixed overall action convention,
\[
 -\Re S^{(2)}(\xi)
 =\frac14\sum_{a<b}k_{ab}
   \bigl|(\xi_b-\xi_a)\times m_{ab}\bigr|^2 .
\]
Its kernel is exactly the common $\SUtwo$ gauge
$\xi_0=\cdots=\xi_4$. Hence after fixing one group element the real Hessian is
strictly negative definite.
\end{proposition}

\begin{proof}
For a normalized spinor $|m\rangle$ and a small relative rotation
$g=\exp(i\,d\!\cdot\!\sigma/2)$,
\[
 |\langle m|g|m\rangle|^2
 =1-\frac14\bigl(|d|^2-(m\!\cdot d)^2\bigr)+O(|d|^3),
\]
so each term $2k_{ab}\log\langle m_{ab}|X_a^{-1}X_b|m_{ab}\rangle$
contributes
$-\tfrac14 k_{ab}|(\xi_b-\xi_a)\times m_{ab}|^2$ to the real quadratic
form. Thus a zero mode must satisfy
$\xi_b-\xi_a\parallel m_{ab}$ on every edge.

Gauge-fix $\xi_0=0$. Exact symbolic facet/bivector arithmetic gives, for both
chiralities,
\[
\begin{array}{lll}
\text{regular:}&
 \det(b_{01},b_{02},b_{12})=\pm\frac{25}{95551488},&
 |b_{03}\!\times b_{13}|^2=|b_{04}\!\times b_{14}|^2
 =\frac{25}{13759414272},\\[3pt]
\text{DK:}&
 \det(b_{01},b_{02},b_{12})=\pm\frac{1}{5971968},&
 |b_{03}\!\times b_{13}|^2=\frac{1}{859963392},\quad
 |b_{04}\!\times b_{14}|^2=\frac{1}{429981696}.
\end{array}
\]
Normalization of the nonzero $b_{ab}$ does not change these independence
statements. The first determinant and the edge conditions for $01,02,12$
force $\xi_1=\xi_2=0$. The two nonparallel-pair identities then force
$\xi_3=0$ and $\xi_4=0$. Therefore the sum of squares has no kernel beyond
the common gauge. Since all $k_{ab}>0$, the gauge-fixed real Hessian is
strictly negative definite. The exact certificate is reproduced by
\texttt{w2\_rehessian\_kernel\_exact.py}; its frozen output is
\texttt{W2-REHESSIAN-KERNEL-EXACT.log}.
\end{proof}

\begin{proposition}\label{prop:hess}
On the regular and on the DK simplex, each chiral $15\times15$ complex Hessian
has exactly three null directions, the common $\SUtwo$ gauge, and its
gauge-fixed $12\times12$ block is nonsingular. Consequently the full doubled
$30\times30$ Hessian has exactly six $\SpinFour$ gauge null directions and the
gauge-fixed $24\times24$ Hessian is nonsingular.
\end{proposition}

\begin{proof}
Write a gauge-fixed chiral Hessian as $H=A+iB$, where
$A=\Re H$ and $A,B$ are real symmetric matrices. Proposition~\ref{prop:redef}
gives $A<0$. If $Hz=0$ for a nonzero complex vector $z$, then
\[
 0=\Re(z^\dagger Hz)=z^\dagger A z<0,
\]
a contradiction. Thus the gauge-fixed chiral Hessian is invertible. Before
gauge fixing, the common left $\SUtwo$ action supplies the three and only three
null directions. The doubled Hessian is block diagonal in the two
chiralities, so the nullities add.
\end{proof}

The earlier numerical calculations are retained as independent checks, not as
the proof. Two implementations sharing no Hessian code agree on the regular
and DK spectra up to their known normalization factors; finite-difference
Richardson extrapolation is stable; and a separately constructed doubled
matrix has six gauge nulls. On ten additional random nondegenerate
four-simplices the same chiral kernel pattern is observed. Kami\'nski and
Sahlmann's analytic nondegeneracy theorem for geometric EPRL stationary points
is consistent with this result~\cite{KaminskiSahlmann2019}, but the present
regular/DK proof does not depend on importing their spinor-variable Hessian.

For reference, the numerical conditioning checks give
\[
\begin{array}{c|cc}
 & \text{regular} & \text{DK}\\ \hline
\max|\mathrm{ev}|/\min_{\ne0}|\mathrm{ev}|,\ 15\times15
 &1.936&3.854\\
\sigma_{\min},\ \text{gauge-fixed }12\times12
 &0.1812404&0.0014178\ /\ 0.0013912 .
\end{array}
\]
These values are evidence of numerical stability only; nondegeneracy itself is
the exact consequence of Proposition~\ref{prop:redef}.

% ===========================================================================
\section{Main theorem}\label{sec:main}

\begin{theorem}\label{thm:main}
Under Assumption~\ref{ass:data}, for the regular or the DK simplex,
\[
\mathcal A_\sigma[\Psi]
=C_\sigma\,\lambda^{-p}\,
\exp\Bigl[i\sigma\,2\sum_f j_f\Theta_f+i\phi_0\Bigr]\bigl(1+O(\lambda^{-1/2+\varepsilon})\bigr),
\qquad C_\sigma\neq0,
\]
where $\Theta_f$ are the dihedral angles, $\phi_0$ is a fixed boundary-convention
phase independent of the bulk, and $p$ is fixed by the $24$-dimensional Gaussian
and the polynomial normalisations. The conjugate orientation contributes to
$\mathcal A_{-\sigma}$, not to $\mathcal A_\sigma$.
\end{theorem}

\begin{proof}
Lemma~\ref{lem:conc} controls the spectral cusp of $|Q|^{1/2}$: the
nonpositive-$Q$ tail of each coherent tetrahedral factor is
$O(\lambda^{-\infty})$. By \eqref{eq:reduce} the volume-weighted integrand is
therefore a finite sum of coherent terms up to $O(\lambda^{-N})$. Theorem~\ref{thm:boundary} bounds
everything outside a neighbourhood of the geometric saddle by $O(\lambda^{-3/2})$
relative to it. At the saddle, Proposition~\ref{prop:hess} gives a nondegenerate
Gaussian and $|\Omega_c|>0$, so standard stationary phase applies. The leading
coefficient is nonzero because the Gaussian is nondegenerate, $|\Omega_c|$ is
positive there, and the volume weights are positive on the retained sign sector.
The phase is that of Section~\ref{sec:critical}.
\end{proof}

The stated relative error $O(\lambda^{-1/2+\varepsilon})$, for every
$\varepsilon>0$, comes from Lemma~\ref{lem:coeff} and needs neither the radial
count of~\cite{W2vertex} nor a symbol expansion. The rank-one boundary term is
smaller, $O(\lambda^{-3/2})$ by the absolute-value argument of
Theorem~\ref{thm:boundary}, so it is the $k=1$ volume-weight term that sets the
rate. A sharper $O(\lambda^{-1})$ would require the symbol expansion at the
saddle, which is not carried out here. Nothing here uses the finite-spin observation that the
wrong-sign weight of the model's own state is well fitted by $e^{-1.04-0.99j}$ for
$j\le19$. Its windowed rates drift, and we report the fit as an observation only.

% ===========================================================================
\section{Gluing and the area-Regge action}\label{sec:r3}

\textbf{This section is formal calculus of the vertex phase extended to
independent area variables. It is not a stationary equation derived for the
coherently glued W2 amplitude}, for the reason given at the end of the section.

Glue W2 vertices on an oriented simplicial complex with the full intertwiner
identity on shared tetrahedra and coherent boundary data compatible across each
shared tetrahedron. Orientation compatibility fixes the relative $\sigma_v$. By
Theorem~\ref{thm:main}, a geometric history carries the phase
$\sum_v\sigma_v\sum_{f\subset v}A_f\Theta_{vf}$ with $A_f=2j_f$.

\begin{proposition}[Area-Regge stationarity]\label{prop:r3}
With $\delta_f=2\pi-\sum_{v\ni f}\Theta_{vf}$, the Schläfli identity
$\sum_{f\subset v}A_f\,d\Theta_{vf}=0$ gives, for internal faces,
$dS_{\mathrm{int}}=\sum_f\delta_f\,dA_f$ up to terms $2\pi A_f$, which change no
amplitude because $A_f=2j_f\in\mathbb Z$. The stationary condition in an internal area
is $\delta_f\in2\pi\mathbb Z$. In a neighbourhood of a flat background small
enough that $|\delta_f|<\pi$ for every internal face, this forces $\delta_f=0$.
\end{proposition}

For the discrete spin sum, Poisson resummation gives the same condition:
constructive interference requires $\partial S/\partial A_f\in2\pi\mathbb Z$,
and the principal branch is $\delta_f=0$. This derivation was verified blind
against acceptance criteria committed before the derivation was read
(Appendix~\ref{app:repro}).

This is the stationary equation of the \emph{area}-Regge action. It is not length
Regge, and we do not promote it to length Regge.

\paragraph{Why this is not yet a statement about glued W2.}
Varying every internal $A_f$ independently is not available on the sector
coherent gluing actually retains. Once the shapes on a shared tetrahedron match
(Section~\ref{sec:r4}), the areas lie on the length-geometric submanifold
$A_f=A_f(l_e)$, and
\[
dS=\sum_f\delta_f\,dA_f
 =\sum_e\Bigl(\sum_f\delta_f\frac{\partial A_f}{\partial l_e}\Bigr)dl_e,
\]
so stationarity gives $\sum_f(\partial A_f/\partial l_e)\,\delta_f=0$, the
\emph{length}-Regge condition, and not $\delta_f=0$ face by face. Independent
$dA_f$, and with them the area-Regge equations, require exactly the
shape-mismatch directions that coherent gluing suppresses. Proposition~\ref{prop:r3}
therefore stands as calculus on independent area variables, not as a derived
consequence of the glued amplitude.

% ===========================================================================
\section{The Dittrich--Kogios limit}\label{sec:r4}

Dittrich and Kogios analyse the area-Regge action on a regular, centrally
subdivided hypercubical lattice~\cite{DittrichKogios2023}. They show that its
continuum limit reproduces at leading order the graviton dynamics of general
relativity, with the first correction at second order in the lattice constant
and quadratic in the Weyl curvature, and they note that this holds contrary to
expectations from the flatness problem. The condition $\delta_f=0$ of
Proposition~\ref{prop:r3} therefore does not by itself preclude graviton dynamics
under refinement.

Their result is obtained from a background-field expansion about flat space on
that one lattice family, and they leave the additional suppression terms of
effective spin-foam models on that triangulation to future work. We would in any
case import it only at that scope, and claim no continuum theorem for a different
triangulation, for a nonperturbative model, or for any other implementation of
the simplicity constraints. As the next subsection explains, we do not reach even
that scope.

\subsection*{The gluing obstruction}

An earlier version of this paper asserted a corollary here: that coherently glued
W2 histories on the DK lattice carry the area-Regge action, so that the
Dittrich--Kogios theorem supplies linearised graviton dynamics. \textbf{That
inference does not follow, and we withdraw it.} The obstruction was identified in
hostile review~\cite{PR357} and is the shape-matching problem in its standard
form~\cite{DittrichSpeziale2008}.

Contract two adjacent vertices through the intertwiner identity on a shared
tetrahedron. If the asymptotics of both vertices are to be supplied by
Theorem~\ref{thm:main}, both sides must present coherent intertwiners
$\iota_L,\iota_R$, and the contraction produces their overlap
\[
\langle\iota_L|\iota_R\rangle
\propto\int dg\prod_f\langle j_f,n^L_f|g|j_f,n^R_f\rangle .
\]
Every factor has modulus at most one, and the real exponent vanishes only if a
single rotation carries all four normals on one side to those on the other, i.e.\
only if the two tetrahedral shapes agree. Otherwise compactness gives a strictly
negative maximum and the overlap is exponentially suppressed, while the
normalisations are only polynomial.

This is exactly the freedom Dittrich and Kogios rely on. Their configuration
space contains neighbouring four-simplices inducing \emph{different} shapes on
their shared tetrahedron, and the resulting shape-mismatch, area-metric degrees of
freedom are integrated out in their derivation. So the chain faces a fork:

\begin{itemize}
\item keep the internal data in the coherent regime of Theorem~\ref{thm:main},
and the shared shapes match, so the DK configuration space has \emph{not} been
obtained; or
\item let the full intertwiner sum explore non-shape-matched data, and
Theorem~\ref{thm:main} supplies no asymptotics for those components.
\end{itemize}

\paragraph{The obstruction fixes a scale, it does not forbid the modes.}
The suppression above is stated at \emph{fixed} mismatch. Writing
$|\langle\iota_L|\iota_R\rangle|\sim e^{\lambda E(\delta)}$ for a one-parameter
shape deformation $\delta$, we find numerically that $E$ is quadratic to leading order,
$E(\delta)=-c\,\delta^2+O(\delta^3)$ with $c>0$ fixed by the areas, so at the
fluctuation scale
$\delta=\epsilon/\sqrt\lambda$ the total exponent $\lambda E=-c\,\epsilon^2$ is
\emph{independent of $\lambda$}. Measured over $\lambda=100,400,1600$ the value
is constant to three figures at each $\epsilon$
(\texttt{shape\_mismatch\_scale.py}). Shape mismatches of order $\lambda^{-1/2}$
are therefore not \emph{exponentially} suppressed in $\lambda$, which is the scale
at which effective spin-foam models carry their area-metric degrees of freedom.
They remain finitely suppressed as a function of $\epsilon$, and integrating those
variables can still generate powers of $\lambda^{-1/2}$; a single overlap does not
fix their net weight in the glued amplitude.

That does not restore the corollary. It says only that the obstruction bounds the
\emph{size} of the surviving mismatch rather than excluding it, so the route must
be taken at the fluctuation scale rather than at fixed shape mismatch. Closing it
still requires deriving the effective glued action after the shared intertwiner
contraction, including the measure over the surviving mismatch modes and the
stationary-phase prefactor, and showing that its quadratic Hessian on the DK
lattice has the required continuum limit. That is not done here, and until it is,
this paper proves a vertex theorem and nothing about the continuum. The remaining
conditions noted in earlier drafts stand as well: Theorem~\ref{thm:main} does not
control the bulk spin sum, and a history must lie in the support of the state
sum.

% ===========================================================================
\section{The model's own tetrahedral state}\label{sec:selected}

The main theorem uses coherent boundary data. The model also selects its own
tetrahedral state, the ground state $\psi$ of $H_+$, and it is natural to ask
whether Theorem~\ref{thm:main} holds with $\iota_e$ replaced by $\psi$. This is
not implied by Lemma~\ref{lem:conc}: $\psi$ is not a coherent intertwiner, and the
two states differ measurably (at $j=8$ the wrong-sign weight is
$1.5\times10^{-4}$ for $\psi$ and $1.9\times10^{-7}$ for $\iota$). Two independent
reviews found that conflating them is the easiest error to make in this chain.

\paragraph{Exact reduction.}
Let $s_1$ be the smallest singular value of $Z$, so $ZZ^\dagger\psi=s_1^2\psi$. By
\eqref{eq:QZ},
\[
Q\psi=\frac{Z^\dagger Z\psi-s_1^2\psi}{3\sqrt3},\qquad
\operatorname{Var}_\psi Q=\tfrac1{27}\operatorname{Var}_\psi(Z^\dagger Z).
\]
When $2j+1\not\equiv0\pmod3$, $s_1=0$ and Chebyshev gives exactly
\[
p_{\le0}(\psi)\le\frac{\operatorname{Var}_\psi Q}{\langle Q\rangle_\psi^2}
=\frac{\operatorname{Var}_\psi(Z^\dagger Z)}{\langle Z^\dagger Z\rangle_\psi^2}.
\]

\begin{conjecture}\label{conj:eig}
The selected state is an asymptotic eigenvector of $Z^\dagger Z$ in relative
variance: $\operatorname{Var}_\psi(Z^\dagger Z)=o(\langle Z^\dagger Z\rangle_\psi^2)$
as $j\to\infty$.
\end{conjecture}

Only $o(1)$ is needed at leading order. Numerically, for $j\le20$, the reduction
holds to relative precision $10^{-12}$, the Chebyshev ratio falls approximately
as $1.5\,j^{-3}$, $\langle Q\rangle/j^3\to0.80$, and the bound lies above the
directly computed wrong-sign weight at every spin. A crude operator-norm estimate
loses by $j^4$, so a proof must use the structure of $\psi$.

\paragraph{Uniqueness.}
For $d=2j+1\le41$: $s_1>0$ if and only if $d\equiv0\pmod3$; the smallest singular
value is simple at every $d$, with gap between $1.5$ and $182$; and the nullity is
exactly one when $d\not\equiv0\pmod3$. The pattern suggests a $\mathbb Z_3$
mechanism from the cyclic relabelling $a\to b\to c$, under which
$Z\mapsto\omega^{\pm1}Z$. We have not proved it. Since $s_1$ decays roughly by a
factor $0.3$ per step of three in $d$, reaching $3\times10^{-6}$ at $d=39$,
numerics cannot settle uniqueness much beyond $j\approx25$.

% ===========================================================================
\section{Claim boundary}\label{sec:claims}

\begin{itemize}
\item \textbf{Proved.} Theorem~\ref{thm:main} for coherent boundary data of the
regular and DK simplices: the uniform coherent-intertwiner norm is proved in
Lemma~\ref{lem:unif}, and Propositions~\ref{prop:hess} and~\ref{prop:redef}
are certified by the exact real-Hessian kernel argument of
Section~\ref{sec:s4};
Proposition~\ref{prop:r3} \emph{only} as formal calculus on independent area
variables; it is \textbf{not} established for the coherently glued amplitude,
where shape matching leaves the length-Regge condition instead
(Section~\ref{sec:r3}).
\item \textbf{Withdrawn.} The continuum inference. Gluing coherent vertices
forces matched shapes on shared tetrahedra, which is the freedom the
Dittrich--Kogios configuration space requires, so their theorem cannot be applied
to glued W2 histories as this paper constructs them (Section~\ref{sec:r4}).
\item \textbf{Conjectured.} Conjecture~\ref{conj:eig}; uniqueness of the selected
state for all $j$.
\item \textbf{Not claimed.} Lorentzian signature; nonlinear curvature; control of
the bulk spin sum of glued histories; dominance for boundary data other than
coherent data of nondegenerate simplices (the Hessian is checked explicitly on
two); a continuum theorem for generic triangulations; any exponential rate as a
theorem; a derivation of general relativity.
\end{itemize}

% ===========================================================================
\section{Relation to existing work}\label{sec:related}

The large-spin analysis of four-simplex amplitudes through coherent states and a
doubled chiral action is due to Barrett et al.~\cite{BarrettEtAl2009}, whose
critical-point classification we use directly; the SU(2) $15j$ case
is~\cite{BFH2009}. Baez, Christensen and Egan found numerical evidence that the Riemannian $10j$ is
dominated by degenerate configurations~\cite{BCE2002}, and Barrett and Steele's
stationary-phase analysis accounts for it; the exceptional families that motivated
the boundary analysis of Section~\ref{sec:boundary} are
Barrett--Steele's~\cite{BarrettSteele2003}. Coherent intertwiners are
Livine--Speziale's~\cite{LivineSpeziale2007}. The appearance of both orientations,
as a cosine of the Regge action, and its cure by a projector selecting the
Einstein--Hilbert sector are the subject of Engle's proper vertex~\cite{Engle2011}.
W2 is a close relative, differing in that the amplitude analysed here fixes
orientation by the integration domain and uses no additional EH boundary
projector. Causal restrictions go
back to Livine and Oriti~\cite{LivineOriti2002}. The canonical measure is
Chaharsough Shirazi and Engle's~\cite{ShiraziEngle2014}. The area-Regge route to
gravitons is Dittrich and Kogios's~\cite{DittrichKogios2023}, and area-angle
variables are due to Dittrich and Speziale~\cite{DittrichSpeziale2008}. The algebra of
the four-valent volume operator is the Racah algebra~\cite{Racah2019}, and the
quantum tetrahedron and the large-spin behaviour of its volume are due to Baez
and Barrett~\cite{BaezBarrett1999} and Schliemann~\cite{Schliemann2013}.

% ===========================================================================
\section{Discussion}\label{sec:disc}

The vertex does what an orientation-resolved vertex should. On nondegenerate
coherent data it has one geometric saddle, with one orientation, and its phase is
the area-Regge action. The degenerate configurations that dominate the
Barrett--Crane $10j$ are present in W2 as the rank-one orbits, and they lose to
the geometric saddle by a power of $\lambda$ because the canonical measure
vanishes cubically on them.

This is a vertex theorem. What stands between it and a derivation of gravity is
what stands in front of every spin-foam model: control of the bulk spin sum on
glued histories, where the flatness problem lives; nonlinear curvature; and
Lorentzian signature. The Dittrich--Kogios result shows that the area-Regge route
is not closed by flatness at the linearised level. Whether it survives beyond that
is not addressed here.

% ===========================================================================
\appendix
\section{Reproducibility and review trail}\label{app:repro}

\paragraph{Frozen objects.}
The W2 stationary equations, Hessian and sign-concentration argument under review
are frozen at commit \texttt{818822c841} on branch
\texttt{sol/w2-epsilon-restriction-audit-2026-09-18}. Independent evidence is on
branch \texttt{worktree-astra-r2-adjudication-2026-09-19} (PR~\#350):
\texttt{w2\_hessian.py} and \texttt{w2\_hessian\_dk.py} (Section~\ref{sec:s4});
\texttt{nongeometric\_criticals.py} (Proposition~\ref{prop:crit});
\texttt{boundary\_strata\_res.py} (Table~\ref{tab:strata});
\texttt{re\_hessian\_definite.py} (numerical conditioning cross-checks);
\texttt{xdx\_identity.py}, \texttt{b1\_moments.py} and \texttt{b4\_mod3.py}
(Section~\ref{sec:selected}); \texttt{R3-PREREGISTRATION.md} and
\texttt{R3-VERIFICATION.md} (Section~\ref{sec:r3}). The exact Hessian
certificate is independently frozen on branch
\texttt{sol/w2-boundary-dominance-2026-09-19} as
\texttt{w2\_rehessian\_kernel\_exact.py} (SHA-256
\texttt{7935701ea3228e4649660ac40e2b724b64feebd3af6a443f06093f6b08fc472c})
and \texttt{W2-REHESSIAN-KERNEL-EXACT.log} (SHA-256
\texttt{a806adb55d2eb9e5f87117e69a720e066fde4cbfb4058c39cf8a661ebf651ffe}).

\paragraph{Reviews that changed this text.}
\begin{itemize}
\item PR~\#355, hostile review of S4/S5: established the analytic Hessian, the
factor $72$ and the six nulls; identified the conflation of coherent and
selected states (Section~\ref{sec:selected}); raised the boundary objection
answered by Theorem~\ref{thm:boundary}; required the admissible sequence
$j_f=\lambda k_f+O(1)$.
\item A context-free review of the sign-concentration argument: the
zero-eigenspace correction ($p_{\le0}$ rather than $p_-$), polynomial
absorption of the super-polynomial tail against the volume insertions, and the
requirement that the moment statement hold for each fixed $n$.
\item A context-free review commissioned on the frozen commit returned
UNRESOLVED because it could not obtain the manuscript, and reviewed the brief
instead. Two of its findings applied regardless and are incorporated: the
Livine--Speziale norm is an asymptotic evaluation rather than a stated uniform
inequality, which is now said explicitly in Lemma~\ref{lem:conc} together with
the finiteness that makes it sufficient; and the Dittrich--Kogios import must
name their background-field expansion and their own deferral of effective
spin-foam suppression terms, which Section~\ref{sec:r4} now does.
\item A second context-free review of the repaired manuscript~\cite{PR358}
returned ACCEPT\_WITH\_SCOPE\_HARDENING and produced the present version. It
found the $k=1$ case of Lemma~\ref{lem:coeff} invalid --- a scalar first moment
had been substituted for a vector norm --- and supplied the interpolation
argument that replaces it, which also yields the $O(\lambda^{-1/2+\varepsilon})$
remainder without appealing to the unpublished companion. It found that the
continuum withdrawal had not been propagated into Section~\ref{sec:r3}, where
independently varied areas are unavailable on the shape-matched sector; that
section is now labelled as formal calculus, with the length-Regge condition
recorded as what glued W2 actually gives. It hardened the shape-mismatch scale
claim from ``not suppressed'' to ``not exponentially suppressed'', sharpened the
uniformity discussion in Lemma~\ref{lem:conc} around the $O(1)$ closure defect,
required the boundary region in Theorem~\ref{thm:boundary} to keep a positive
distance from the rank-one orbits, and noted that the evidence bundle claimed a
completeness it did not have. It checked and cleared the rank-one absolute-value
estimate, the cubic cofactor count, the separation of boundary-state from
$X$-space statements, and the repaired citations.
\item The R4 hostile review~\cite{R4review}: placement of the
degenerate-dominance condition, which Theorem~\ref{thm:boundary} now discharges at
the vertex, and the admissibility and area-independence conditions of
Section~\ref{sec:r4}.
\item A context-free review of the frozen manuscript~\cite{PR357} returned
REJECT and is the reason for the present version. It identified the gluing
obstruction of Section~\ref{sec:r4}, which removed the continuum corollary and
the former subtitle; the $2j_f+1$ phase defect, repaired by using $A_f=2j_f$; the
insufficiency of the earlier finiteness argument in Lemma~\ref{lem:conc}; the
unproved coefficient ratios behind the $O(\lambda^{-1})$ remainder; the central
sign lifts in Proposition~\ref{prop:crit}; the evidentiary status of
Propositions~\ref{prop:hess} and~\ref{prop:redef}; a spliced bibliography entry
for the Racah algebra; and an overstatement of what Baez, Christensen and Egan
proved. It also checked and cleared the Taylor remainder of~\eqref{eq:reduce},
the direction of the orthogonal-projection inequality used in
Lemma~\ref{lem:conc}, and the naming of the numerical conditioning quantities
retained in Section~\ref{sec:s4}.
\item An earlier hostile-review brief withdrew two of its own claims: that
$O(\lambda^{-\infty})$ requires constants uniform in $n$ (false), and that the
chiral blocks might couple (excluded by \eqref{eq:action}).
\item A later independent boundary-dominance lane supplied an exact
sum-of-squares/kernel certificate for the real chiral Hessian on both the
regular and DK backgrounds. Section~\ref{sec:s4} now uses that certificate to
remove the last floating-point nondegeneracy hypothesis: strict negativity of
$\Re H$ modulo gauge implies invertibility of the complex gauge-fixed Hessian.
\end{itemize}

\paragraph{Retracted claims.}
Five earlier claims of this programme were retracted, including a sign-structure
claim (PR~\#351). The measured $e^{-0.995j}$ law is not claimed.

\begin{thebibliography}{99}
\bibitem{BarrettEtAl2009} J.~W. Barrett, R.~J. Dowdall, W.~J. Fairbairn,
H.~Gomes and F.~Hellmann, \emph{Asymptotic analysis of the EPRL four-simplex
amplitude}, J.\ Math.\ Phys.\ \textbf{50} (2009) 112504.
\bibitem{BFH2009} J.~W. Barrett, W.~J. Fairbairn and F.~Hellmann,
\emph{Quantum gravity asymptotics from the $SU(2)$ $15j$ symbol},
Int.\ J.\ Mod.\ Phys.\ A \textbf{25} (2010) 2897.
\bibitem{BCE2002} J.~C. Baez, J.~D. Christensen and G.~Egan,
\emph{Asymptotics of $10j$ symbols}, Class.\ Quant.\ Grav.\ \textbf{19} (2002) 6489.
\bibitem{BarrettSteele2003} J.~W. Barrett and C.~M. Steele,
\emph{Asymptotics of relativistic spin networks},
Class.\ Quant.\ Grav.\ \textbf{20} (2003) 1341.
\bibitem{DittrichKogios2023} B.~Dittrich and A.~Kogios,
\emph{From spin foams to area metric dynamics to gravitons},
Class.\ Quant.\ Grav.\ \textbf{40} (2023) 095011, arXiv:2203.02409.
\bibitem{DittrichSpeziale2008} B.~Dittrich and S.~Speziale,
\emph{Area-angle variables for general relativity},
New J.\ Phys.\ \textbf{10} (2008) 083006.
\bibitem{Engle2011} J.~Engle,
\emph{A proposed proper EPRL vertex amplitude},
Phys.\ Rev.\ D \textbf{87} (2013) 084048.
\bibitem{LivineOriti2002} E.~R. Livine and D.~Oriti,
\emph{Implementing causality in the spin foam quantum geometry},
Nucl.\ Phys.\ B \textbf{663} (2003) 231.
\bibitem{LivineSpeziale2007} E.~R. Livine and S.~Speziale,
\emph{A new spinfoam vertex for quantum gravity},
Phys.\ Rev.\ D \textbf{76} (2007) 084028.
\bibitem{Racah2019} N.~Crampé, L.~Poulain d'Andecy and L.~Vinet,
\emph{Temperley--Lieb, Brauer and Racah algebras and other centralizers of
$\mathfrak{su}(2)$}, arXiv:1905.06346.
\bibitem{KaminskiSahlmann2019} W.~Kamiński and H.~Sahlmann,
\emph{The Hessian in spin foam models},
Ann.\ Henri Poincaré \textbf{20} (2019) 3927, arXiv:1906.05258.
\bibitem{ConradyFreidel2008} F.~Conrady and L.~Freidel,
\emph{On the semiclassical limit of 4d spin foam models},
Phys.\ Rev.\ D \textbf{78} (2008) 104023.
\bibitem{BaezBarrett1999} J.~C. Baez and J.~W. Barrett,
\emph{The quantum tetrahedron in $3$ and $4$ dimensions},
Adv.\ Theor.\ Math.\ Phys.\ \textbf{3} (1999) 815.
\bibitem{Schliemann2013} J.~Schliemann,
\emph{The large-volume limit of a quantum tetrahedron is a quantum harmonic
oscillator}, Class.\ Quant.\ Grav.\ \textbf{30} (2013) 235018, arXiv:1307.5979.
\bibitem{Regge1961} T.~Regge,
\emph{General relativity without coordinates},
Nuovo Cimento \textbf{19} (1961) 558.
\bibitem{ShiraziEngle2014} A.~Chaharsough Shirazi and J.~Engle,
\emph{Purely geometric path integral for spin foams},
Class.\ Quant.\ Grav.\ \textbf{31} (2014) 075010.
\bibitem{W2vertex} T.~Croydon-McRae,
\emph{An oriented finite-spin vertex with retained tetrahedral geometry:
construction, canonical measure, and the gluing obstruction},
preprint v0.2 (2026), \texttt{doi:10.6084/m9.figshare.33936508}.
\bibitem{PR355} Hostile review of W2 S4/S5, physics-lane PR~\#355, branch
\texttt{claude/review-w2-s4-s5-hostile-2026-09-19}, commit \texttt{92aea65855}.
\bibitem{PR357} Context-free hostile review of this manuscript at commit
\texttt{c2d47cfc}, 19 September 2026; disposition REJECT.
\bibitem{PR358} Second context-free hostile review of this manuscript,
19 September 2026; disposition ACCEPT\_WITH\_SCOPE\_HARDENING.
\bibitem{R4review} Hostile review of R4,
\texttt{orthogonal-validation/astra-r2-adjudication-2026-09-19/R4-HOSTILE-REVIEW.md}.
\end{thebibliography}

\end{document}
